% libre-compendium
%
% Copyright (C) 2026 Richard Qin
%
% This document is dual-licensed under either:
% - GNU Free Documentation License, Version 1.3 or any later version
%   published by the Free Software Foundation; with no Invariant Sections,
%   no Front-Cover Texts, and no Back-Cover Texts (see COPYING.DOC); or
% - Creative Commons Attribution-ShareAlike 4.0 International License
%   (see LICENSE-CC-BY-SA-4.0).

% Exercises migrated from problems01.tex and solutions01.tex.
\begin{exercise}[level=basic]
写出二次型
\[
Q=4x_1^2+x_2^2-3x_3^2+6x_1x_2-2x_1x_3+8x_2x_3
\]
的对称矩阵。
\begin{solution}
对应矩阵为
\[
\begin{pmatrix}4&3&-1\\3&1&4\\-1&4&-3\end{pmatrix}.
\]
交叉项系数依次折半为 $3,-1,4$。
\end{solution}
\end{exercise}

\begin{exercise}[level=basic]
设 $A=\begin{pmatrix}1&2\\2&-1\end{pmatrix}$。写出 $Q(x)=x^{\mathsf T}Ax$ 及 $B(x,y)=x^{\mathsf T}Ay$。
\begin{solution}
\[
Q=x_1^2+4x_1x_2-x_2^2,
\]
\[
B(x,y)=x_1y_1+2x_1y_2+2x_2y_1-x_2y_2.
\]
\end{solution}
\end{exercise}

\begin{exercise}[level=basic]
证明对任意实反对称矩阵 $K$，恒有 $x^{\mathsf T}Kx=0$。
\begin{solution}
$x^{\mathsf T}Kx$ 是 $1\times1$ 矩阵，故等于其转置；另一方面
\[
(x^{\mathsf T}Kx)^{\mathsf T}=x^{\mathsf T}K^{\mathsf T}x=-x^{\mathsf T}Kx.
\]
因此它等于自身的相反数，只能为零。
\end{solution}
\end{exercise}

\begin{exercise}[level=basic]
判断下列映射是否为双线性型；若是，再判断是否对称：
\begin{enumerate}[label=(\arabic*)]
\item $B(x,y)=x_1y_1+x_1y_2$，$x,y\in\R^2$；
\item $B(x,y)=x_1y_1+3$；
\item $B(X,Y)=\tr(XY)$，$X,Y\in\R^{n\times n}$。
\end{enumerate}
\begin{solution}
(1) 是双线性型，矩阵为 $\begin{pmatrix}1&1\\0&0\end{pmatrix}$，不对称。(2) 不是双线性型，因为 $B(0,y)=3\ne0$。(3) 是双线性型，并且
\[
\tr(XY)=\tr(YX),
\]
故对称。这里使用的是方阵乘积迹的循环性质。
\end{solution}
\end{exercise}

\begin{exercise}[level=basic]
用极化公式由 $Q(x_1,x_2)=x_1^2+4x_1x_2+7x_2^2$ 求对应的对称双线性型。
\begin{solution}
矩阵为 $A=\begin{pmatrix}1&2\\2&7\end{pmatrix}$，故
\[
B(x,y)=x_1y_1+2x_1y_2+2x_2y_1+7x_2y_2.
\]
也可把 $x+y$ 代入 $Q$ 后按 $[Q(x+y)-Q(x)-Q(y)]/2$ 展开。
\end{solution}
\end{exercise}

\begin{exercise}[level=basic]
证明矩阵相合是等价关系。
\begin{solution}
自反性：$A=I^{\mathsf T}AI$。对称性：若 $B=P^{\mathsf T}AP$，则 $A=(P^{-1})^{\mathsf T}BP^{-1}$。传递性：若另有 $C=Q^{\mathsf T}BQ$，则 $C=(PQ)^{\mathsf T}A(PQ)$。
\end{solution}
\end{exercise}

\begin{exercise}[level=basic]
将 $Q=3x_1^2+6x_1x_2+2x_2^2$ 用配方法化为对角形，并求正、负惯性指数。
\begin{solution}
\[
Q=3(x_1+x_2)^2-x_2^2.
\]
令 $y_1=x_1+x_2,y_2=x_2$，得 $Q=3y_1^2-y_2^2$，故正、负惯性指数均为 $1$。
\end{solution}
\end{exercise}

\begin{exercise}[level=basic]
判断 $\diag(2,0,3)$、$\diag(-1,-2)$、$\diag(1,-1,0)$ 分别属于正定、半正定、负定、半负定或不定中的哪一类。
\begin{solution}
$\diag(2,0,3)$ 半正定但不正定；$\diag(-1,-2)$ 负定；$\diag(1,-1,0)$ 不定。
\end{solution}
\end{exercise}

\begin{exercise}[level=standard]
将
\[
Q=2x_1x_2+2x_1x_3+2x_2x_3
\]
化为规范形，并求秩和惯性指数。
\begin{solution}
令 $u=x_1+x_2$、$v=x_1-x_2$，则 $2x_1x_2=(u^2-v^2)/2$，而 $x_1+x_2=u$，故
\[
Q=\frac12u^2-\frac12v^2+2ux_3
=\frac12(u+2x_3)^2-\frac12v^2-2x_3^2.
\]
三个非零平方项中一个为正、两个为负。规范形为 $z_1^2-z_2^2-z_3^2$，故 $\rank A=3$，$(p,q)=(1,2)$。
\end{solution}
\end{exercise}

\begin{exercise}[level=standard]
求实数 $a$ 的范围，使
\[
A=\begin{pmatrix}1&a&a\\a&1&a\\a&a&1\end{pmatrix}
\]
正定。
\begin{solution}
顺序主子式为
\[
\Delta_1=1,\qquad \Delta_2=1-a^2,\qquad
\Delta_3=(1-a)^2(1+2a).
\]
要求 $|a|<1$ 且 $1+2a>0$；在 $|a|<1$ 时 $(1-a)^2>0$。故
\[
-\frac12<a<1.
\]
\end{solution}
\end{exercise}

\begin{exercise}[level=standard]
设 $A$ 为实对称矩阵。证明若 $A>0$，则 $A^{-1}>0$。
\begin{solution}
对任意 $x\ne0$，令 $y=A^{-1}x\ne0$，则
\[
x^{\mathsf T}A^{-1}x=(Ay)^{\mathsf T}A^{-1}(Ay)=y^{\mathsf T}Ay>0,
\]
其中使用了 $A^{\mathsf T}=A$。故 $A^{-1}>0$。
\end{solution}
\end{exercise}

\begin{exercise}[level=standard]
设 $A>0$，$B>0$。判断 $AB$ 是否一定为对称正定矩阵；同时证明 $AB$ 的全部特征值均为正实数。
\begin{solution}
$AB$ 不一定对称，因为 $(AB)^{\mathsf T}=BA$，只有 $AB=BA$ 时才对称。由于 $A>0$，存在正定平方根 $A^{1/2}$。矩阵
\[
AB=A^{1/2}(A^{1/2}BA^{1/2})A^{-1/2}
\]
与对称正定矩阵 $A^{1/2}BA^{1/2}$ 相似，故 $AB$ 可对角化且所有特征值均为正实数。
\end{solution}
\end{exercise}

\begin{exercise}[level=standard]
在多项式空间 $P_2(\R)$ 上定义
\[
B(f,g)=f(0)g(0)+f(1)g(1)+f(-1)g(-1).
\]
求其在基 $(1,t,t^2)$ 下的 Gram 矩阵，并判断 $B$ 是否为内积。
\begin{solution}
在 $0,1,-1$ 三点上计算基函数值，得
\[
G=\begin{pmatrix}
3&0&2\\
0&2&0\\
2&0&2
\end{pmatrix}.
\]
其顺序主子式为 $3,6,4$，均为正，故 $G>0$，$B$ 为内积。也可直接说明：若 $B(f,f)=0$，则次数不超过 $2$ 的多项式在三个不同点均为零，只能是零多项式。
\end{solution}
\end{exercise}

\begin{exercise}[level=standard]
设双线性型在基 $\mathcal E$ 下的矩阵为
\[
G=\begin{pmatrix}2&1\\1&2\end{pmatrix},
\]
新基满足 $\widetilde{\mathcal E}=\mathcal E P$，其中
\[
P=\begin{pmatrix}1&1\\0&1\end{pmatrix}.
\]
求新基下的矩阵，并写出旧坐标与新坐标的关系。
\begin{solution}
\[
\widetilde G=P^{\mathsf T}GP
=\begin{pmatrix}2&3\\3&6\end{pmatrix}.
\]
若旧坐标为 $x$、新坐标为 $\widetilde x$，则
\[
x=P\widetilde x,\qquad \widetilde x=P^{-1}x.
\]
\end{solution}
\end{exercise}

\begin{exercise}[level=standard]
设实对称矩阵 $A,B$ 相合。证明 $\rank A=\rank B$，并说明它们的行列式是否一定相等。
\begin{solution}
若 $B=P^{\mathsf T}AP$，则左右乘可逆矩阵不改变秩，故 $\rank B=\rank A$。行列式满足
\[
\det B=(\det P)^2\det A,
\]
因此不一定相等，但其是否为零以及非零时的符号保持不变。
\end{solution}
\end{exercise}

\begin{exercise}[level=standard]
不用求特征值，判断
\[
A=\begin{pmatrix}2&-1&0\\-1&2&-1\\0&-1&2\end{pmatrix}
\]
是否正定。
\begin{solution}
顺序主子式为
\[
\Delta_1=2,\qquad \Delta_2=3,\qquad \Delta_3=4.
\]
全部为正，故 $A$ 正定。
\end{solution}
\end{exercise}

\begin{exercise}[level=advanced]
证明：若二次型 $Q$ 的正惯性指数为 $p$，则满足 $Q|_{W\setminus\{0\}}>0$ 的子空间 $W$ 的最大维数恰为 $p$。
\begin{solution}
惯性定理给出坐标，使
\[
Q=y_1^2+\cdots+y_p^2-y_{p+1}^2-\cdots-y_{p+q}^2.
\]
前 $p$ 个坐标张成的子空间上 $Q$ 正定，故最大维数至少为 $p$。若 $W$ 的维数大于 $p$，从 $W$ 到前 $p$ 个坐标空间的投影必有非零核向量；该向量前 $p$ 个坐标全零，故 $Q\le0$，与 $Q|_{W\setminus\{0\}}>0$ 矛盾。因此最大维数恰为 $p$。
\end{solution}
\end{exercise}

\begin{exercise}[level=advanced]
设 $A$ 为 $m\times n$ 实矩阵。证明 $A^{\mathsf T}A$ 半正定，并证明
\[
\ker(A^{\mathsf T}A)=\ker A,\qquad \rank(A^{\mathsf T}A)=\rank A.
\]
何时 $A^{\mathsf T}A$ 正定？
\begin{solution}
对任意 $x$，
\[
x^{\mathsf T}A^{\mathsf T}Ax=\lVert Ax\rVert^2\ge0,
\]
故半正定。并且
\[
A^{\mathsf T}Ax=0\Rightarrow x^{\mathsf T}A^{\mathsf T}Ax=0
\Rightarrow \lVert Ax\rVert^2=0\Rightarrow Ax=0.
\]
反向包含显然，故两核相等。由秩零度定理可得秩相等。$A^{\mathsf T}A$ 正定当且仅当 $\ker A=\{0\}$，即 $A$ 的列向量线性无关，亦即 $\rank A=n$。
\end{solution}
\end{exercise}

\begin{exercise}[level=advanced]
设 $A=A^{\mathsf T}$。证明存在可逆矩阵 $P$ 使 $P^{\mathsf T}AP=A$ 的矩阵构成一个群。若 $A=I_n$，该群是什么？
\begin{solution}
记
\[
O(A)=\{P\in GL_n(\R):P^{\mathsf T}AP=A\}.
\]
$I\in O(A)$。若 $P,Q\in O(A)$，则
\[
(PQ)^{\mathsf T}A(PQ)=Q^{\mathsf T}(P^{\mathsf T}AP)Q=Q^{\mathsf T}AQ=A.
\]
若 $P\in O(A)$，由 $P^{\mathsf T}AP=A$ 左右乘相应逆矩阵得 $(P^{-1})^{\mathsf T}AP^{-1}=A$。故它对乘法构成群。$A=I_n$ 时得到正交群 $O(n)$。
\end{solution}
\end{exercise}

\begin{exercise}[level=advanced]
设
\[
A(t)=\begin{pmatrix}
1&t&0\\
t&1&t\\
0&t&1
\end{pmatrix}.
\]
随参数 $t$ 的变化，求 $A(t)$ 的正、负、零惯性指数。
\begin{solution}
顺序主子式为
\[
\Delta_1=1,\qquad \Delta_2=1-t^2,\qquad \Delta_3=1-2t^2.
\]
也可直接求特征值 $1,1+\sqrt2t,1-\sqrt2t$。因此：
\begin{enumerate}[label=(\arabic*)]
\item $|t|<1/\sqrt2$ 时，惯性为 $(3,0,0)$；
\item $|t|=1/\sqrt2$ 时，惯性为 $(2,0,1)$；
\item $|t|>1/\sqrt2$ 时，一个特征值为负、两个为正，惯性为 $(2,1,0)$。
\end{enumerate}
\end{solution}
\end{exercise}

\begin{exercise}[level=advanced]
给出一个 $3$ 阶实对称矩阵，使其所有对角元均为正，但矩阵不正定；并解释为什么只检查对角元不能判定正定性。
\begin{solution}
例如
\[
A=\begin{pmatrix}1&2&0\\2&1&0\\0&0&1\end{pmatrix}.
\]
其对角元均为正，但左上 $2$ 阶顺序主子式为 $1-4=-3<0$，故不正定；事实上取 $x=(1,-1,0)^{\mathsf T}$，有 $x^{\mathsf T}Ax=-2<0$。对角元只记录坐标轴方向上的取值，不能控制不同坐标同时非零时交叉项造成的符号变化。
\end{solution}
\end{exercise}

\begin{exercise}[level=advanced]
设 $B$ 是有限维实线性空间上的对称双线性型。证明：$B$ 非退化当且仅当它在任意一组基下的 Gram 矩阵可逆；并说明该性质为何与基无关。
\begin{solution}
定义线性映射 $T_B:V\to V^*$，$T_B(u)=B(u,\cdot)$。$B$ 非退化是指 $\ker T_B=\{0\}$。在任一组基及其对偶基下，$T_B$ 的矩阵正是 Gram 矩阵 $G$，故非退化等价于 $G$ 可逆。换基后 $\widetilde G=P^{\mathsf T}GP$，于是
\[
\det\widetilde G=(\det P)^2\det G.
\]
因 $P$ 可逆，$\det\widetilde G=0$ 当且仅当 $\det G=0$，故该性质与基无关。
\end{solution}
\end{exercise}

\begin{exercise}[level=advanced]
利用 Jacobi 惯性判据求矩阵
\[
A=\begin{pmatrix}
1&2&0\\
2&1&1\\
0&1&-2
\end{pmatrix}
\]
的正、负惯性指数。
\begin{solution}
顺序主子式为
\[
\Delta_0=1,\qquad \Delta_1=1,\qquad
\Delta_2=-3,\qquad \Delta_3=5.
\]
数列 $1,1,-3,5$ 中有一次同号、两次异号，因此正惯性指数为 $1$，负惯性指数为 $2$。
\end{solution}
\end{exercise}

\begin{exercise}[level=advanced]
判断矩阵
\[
A=\begin{pmatrix}
1&a\\
a&b
\end{pmatrix}
\]
半正定的充要条件。
\begin{solution}
二阶实对称矩阵半正定，当且仅当所有主子式非负。因此必须且只需
\[
1\ge0,\qquad b\ge0,\qquad b-a^2\ge0.
\]
第一个条件恒成立，而 $b-a^2\ge0$ 已蕴含 $b\ge0$，故充要条件为
\[
b\ge a^2.
\]
\end{solution}
\end{exercise}

\begin{exercise}[level=advanced]
设 $B>0$，证明分块对称矩阵
\[
\begin{pmatrix}
B&C\\
C^{\mathsf T}&D
\end{pmatrix}
\]
正定，当且仅当 $D-C^{\mathsf T}B^{-1}C>0$。
\begin{solution}
取
\[
P=\begin{pmatrix}I&-B^{-1}C\\0&I\end{pmatrix}.
\]
由于 $B>0$，矩阵 $B$ 可逆。直接计算得
\[
P^{\mathsf T}
\begin{pmatrix}B&C\\C^{\mathsf T}&D\end{pmatrix}
P
=
\begin{pmatrix}
B&0\\
0&D-C^{\mathsf T}B^{-1}C
\end{pmatrix}.
\]
相合保持正定性，而右端分块对角矩阵正定当且仅当两个对角块均正定。已知 $B>0$，故原矩阵正定恰等价于其 Schur 补正定。
\end{solution}
\end{exercise}

\begin{exercise}[level=advanced]
在 $P_2(\R)$ 上定义
\[
B(f,g)=f(0)g(0)+f'(0)g'(0).
\]
求 $B$ 在基 $(1,t,t^2)$ 下的 Gram 矩阵、秩与根空间，并判断是否非退化。
\begin{solution}
设 $f=x_1+x_2t+x_3t^2$、$g=y_1+y_2t+y_3t^2$，则
\[
B(f,g)=x_1y_1+x_2y_2.
\]
故 Gram 矩阵为
\[
G=\diag(1,1,0).
\]
因此 $\rank B=2$。根空间对应 $\ker G$，所以
\[
\operatorname{rad}B=\operatorname{span}\{t^2\}.
\]
根空间非零，故 $B$ 退化。
\end{solution}
\end{exercise}

\begin{exercise}[level=advanced]
设 $A=A^{\mathsf T}\ge0$。证明
\[
\ker A=\ker A^{1/2},
\qquad
\rank A=\rank A^{1/2}.
\]
\begin{solution}
由谱分解，存在正交矩阵 $U$ 使
\[
A=U\diag(\lambda_1,\ldots,\lambda_n)U^{\mathsf T},
\quad
A^{1/2}=U\diag(\sqrt{\lambda_1},\ldots,\sqrt{\lambda_n})U^{\mathsf T}.
\]
$\lambda_i=0$ 当且仅当 $\sqrt{\lambda_i}=0$，所以两个矩阵具有相同的零特征子空间，故核相同。由秩零度定理，它们的秩也相同。
\end{solution}
\end{exercise}

\begin{exercise}[level=advanced]
设
\[
A=\begin{pmatrix}2&1\\1&2\end{pmatrix}.
\]
求 Rayleigh 商的取值范围，并说明何时取得端点。
\begin{solution}
$A$ 的特征值为 $1,3$，对应特征向量可分别取 $(1,-1)^{\mathsf T}$ 与 $(1,1)^{\mathsf T}$。因此
\[
1\le R_A(x)\le3.
\]
当 $x$ 是 $(1,-1)^{\mathsf T}$ 的非零倍数时取最小值；当 $x$ 是 $(1,1)^{\mathsf T}$ 的非零倍数时取最大值。
\end{solution}
\end{exercise}

\begin{exercise}[level=advanced]
在 $\R^3$ 上定义
\[
B(x,y)=x_1y_1+x_2y_2-x_3y_3.
\]
取 $W=\operatorname{span}\{(1,0,1)^{\mathsf T}\}$，求 $W^{\perp_B}$，并判断 $V$ 是否等于 $W\oplus W^{\perp_B}$。
\begin{solution}
设 $v=(a,b,c)^{\mathsf T}$。条件 $v\in W^{\perp_B}$ 等价于
\[
B\bigl(v,(1,0,1)^{\mathsf T}\bigr)=a-c=0.
\]
故
\[
W^{\perp_B}
=\operatorname{span}\{(1,0,1)^{\mathsf T},(0,1,0)^{\mathsf T}\}.
\]
可见 $W\subseteq W^{\perp_B}$，所以交集不为零，不能构成直和。这里
\[
B((1,0,1)^{\mathsf T},(1,0,1)^{\mathsf T})=0,
\]
即 $B$ 在 $W$ 上的限制是退化的。
\end{solution}
\end{exercise}

